\chapter{Conclusion}

This Master's/PFE project developed and studied a reproducible LLM-as-a-Judge reliability laboratory. Its principal result is not that language-model judges are simply reliable or unreliable. Instead, the evidence identifies a profile. On the source-corrected final population, the judges reached 59.20\% agreement with human preference reference labels and showed 96.53\% fixed-condition repeated consistency. At the same time, 17.49\% of decisive paired comparisons changed their mapped outcome after answer-order swapping. These results demonstrate why agreement alone is insufficient for evaluating an automatic judge.

The controlled redundant-length and presentation-format treatments did not obtain stable variant wins in their eligible populations. The source-family experiment did not support a uniform aggregate same-family preference, while revealing strong judge-level heterogeneity. RQ7 showed that mitigation must be presented with both agreement and coverage. DUAL\_SWAP had a small, uncertain agreement change and a material coverage reduction. Strict Multi-Judge Consensus provided a positive secondary improvement relative to its own equal-weight individual-judge comparator on its retained pairs. Their effects are not direct head-to-head measurements.

The project also demonstrates an engineering lesson: final scientific results require an auditable path from exact source text to stored observations, analysis-run identities, release artifacts, API responses, and user interface. The source-text correction, durable terminal-state accounting, API fail-closed rules, and provider-free release validation make the numbers more than dashboard labels. They make the evidence inspectable.

The final recommendation is therefore practical and cautious. Use LLM judges as controlled instruments, not as unquestioned arbiters. Report human-reference agreement together with repeatability, order sensitivity, treatment scope, coverage, and failure states. Keep mitigation comparators explicit. Preserve exact provenance and frozen artifacts. Under these conditions, automated judging can support scalable evaluation while retaining the methodological limits needed for credible research conclusions.
